\section{refined Sp$(2N)$ の表現 [1] と 2 箱}\label{sec:sp}
\cite{AM22} の refined Chern--Simons の定義は、任意のルート系で $k_\alpha = y$ を一様にとり $q_\alpha = q^{(\alpha,\alpha)/2}$ とするもので、C$_N$ では長い根のパラメータが $u = t_{\mathrm{lng}} = t^2$ になる。
\cite{Cherednik11} §Kauffman の C$_n$ の DAHA（$t = t_{\mathrm{sht}}$, $u = t_{\mathrm{lng}}$）を多項式表現の上で実装し、次を確かめた。
% [D11]
\begin{compresult}
\begin{enumerate}
\item C$_2$, C$_3$ で DAHA の関係式が成り立つ。$u = 1$ で D$_N$ の DAHA-Jones の表現 [1] に比 1 で一致する（$N = 2, 3$、$T(3,2)$, $T(5,2)$）。Cherednik の $\mathcal{H}^C_{3,2}(\omega_1; q, t, u, a_{\mathrm{C}})$ に $a_{\mathrm{C}} = -t^{m-1}$ で単項式を除いて一致する（$m = 2, 3, 4$）。
\item $N = 2, \dots, 8$ の値から hyper-polynomial $\mathcal{H}^C_{2m+1,2}(\omega_1)$ を $a_{\mathrm{C}}$（Cherednik の $a$）の多項式として復元した。$T(5,2)$ の式は新しい。
\item SO の規約：$W_{\SO} = $ 鏡像$[\mathcal{H}^C(u{=}1, a_{\mathrm{C}}{=}-\lambda/t^{1/2})]\cdot\dim_V^{\SO}(\lambda)\cdot(t^{1/2}/\lambda)^{2m}$ で、$g^{\mathrm{true}}_{\yone} = W_{\SO} - W^{\cov}_1$ が再現される（一意）。
\item Sp：$\lambda = t^{N+1/2}$（$a = \lambda^2 = t^{2N+1}$。被覆は SO と同じ $\lambda$ の関数で取る）、$a_{\mathrm{C}} = -t^{N-1}$、$u = t^2$。単項式は $t = 1$ の極の打ち消しと $t = q$ での Kauffman の形から SO と同じ $(t^{1/2}/\lambda)^{2m}$ に決まり、$T(2,3)$, $T(2,5)$ で
\[
g^{\mathrm{Sp}}_{\yone}(\ba, \bq, \bt) = -\,g^{\SO}_{\yone}(-\ba, \bq, \bt)
\]
が全係数で成り立つ。
\end{enumerate}
\end{compresult}
$c = 1$ を $\ba$ の偶数冪、$c = 2$ を奇数冪に割り当てているので、これは $C^{\mathrm{Sp}} = -C_{c=1} + C_{c=2}$、すなわち Mariño の「SO と Sp は $c = 1$ の項の符号が違う」が refined でもそのまま成り立つことを意味する。$c$ の割り当てを独立に裏付けるとともに、\cite{AM22} の refined Sp$(2N)$ が開弦でも整合する最初の確認である。

2 箱でも同じ構造を確かめた。C$_N$ の Macdonald 多項式 $P_{2\varepsilon_1}$, $P_{\varepsilon_1+\varepsilon_2}$ を Macdonald 作用素の三角行列から求め、$N = 2, \dots, 9$ の値を $x = t^N$ の有理関数に補間した（$u = 1$ では D$_N$ の 2 箱に一致し、Cherednik の $\mathcal{H}^C_{3,2}(\omega_2)$ に単項式を除いて一致する）。
% [D12][D13][D14]
\begin{compresult}
trefoil の refined Sp$(2N)$ の 2 箱で、
\begin{enumerate}
\item $g^{\mathrm{Sp}}_{\ytwo} + g^{\mathrm{Sp}}_{\yoneone}$ と $(1+t)(g^{\mathrm{Sp}}_{\ytwo} - g^{\mathrm{Sp}}_{\yoneone})$ は整数 Laurent 多項式で、SO と同じ全 $d$ の構造（命題~\ref{prop:trace}）をもつ。
\item $E^{\mathrm{Sp}}(t{=}-1) = E^{\SO}(t{=}-1)$（全係数）。
\item hatted の $p_1^2$ 成分 $(g_{\ytwo} + g_{\yoneone})/z_q$ は Sp$(\ba) = +\,\SO(-\ba)$（表現 [1] は $-$）。$g_R$ そのものは
\begin{align*}
f^{\mathrm{Sp}}_{\ytwo}(\ba) - f^{\SO}_{\ytwo}(-\ba) &= \psi^2[\delta], \qquad f^{\mathrm{Sp}}_{\yoneone}(\ba) - f^{\SO}_{\yoneone}(-\ba) = -\psi^2[\delta], \\
\delta &= \mathrm{even}_{\ba}\big(g^{\SO}_{\yone}\big) + 1 = \frac{(\ba^2 - \bq^2)(\ba^2\bt^2 - 1)(\ba^2\bq^2\bt^2 - 1)}{\ba^6\bq^2\bt^4}
\end{align*}
を満たす（$g$ でも同じ。$\psi^2$ は $q, t, \ba^2$ を 2 乗する Adams 作用素）。
\end{enumerate}
\end{compresult}
冪和の成分では $Z_1^{\mathrm{Sp}}(\ba) = -Z_1^{\SO}(-\ba)$, $Z_{11}^{\mathrm{Sp}}(\ba) = Z_{11}^{\SO}(-\ba)$, $Z_2^{\mathrm{Sp}}(\ba) = Z_2^{\SO}(-\ba) + 1$ で、符号は $(-1)^{|\rho|}$、$p_2$ 方向にだけ定数 $1$ のずれがある。$f$ の差 $\psi^2[\delta]$ のうち $\psi^2[\mathrm{even}(g_{\yone})]$ は、表現 [1] の規則が $\ba$ の偶数部分（$c = 1$）の符号だけを変えるのに対し $d = 2$ の Adams 項 $\tfrac12\psi^2[f_{\yone}]$ の偶数部分がそれに従わないことから生じ、$\psi^2[1]$ は $Z_2$ の定数のずれである。彩色の $W_R$ の段階では $W_R^{\mathrm{Sp}}(\ba) = (-1)^{|R|}\,W_R^{\SO}(-\ba)$ が厳密に成り立ち、冪和への変換係数は $p_2$ の定数項（恒等元のトレース）だけが $-2$ 違う。これが $Z_2$ の定数 $1$ で、unknot の表現 [1] のクロスキャップの $\ba$ 偶数部分（SO で $+1$、Sp で $-1$）に等しい。
2 箱の計算は trefoil に限る（$T(2,5)$ は補間に $N \le 13$ が要り、現在の実装では計算量が大きすぎる）。
$E$ が Sp への射影で符号を変えないことは、unrefined の対応「Sp$(N) = \SO(-N)$ かつ $R \to R^{T}$」で $E_{\ytwo} = -E_{\yoneone}$ を使うと、符号反転と転置が打ち消して $E$ が不変になることと整合する。$p_2$ 方向には表現 [1] の Adams 項 $\psi^2[g_{\yone}]$ が入り、$\psi^2$（$\lambda \to \lambda^2$）は $\ba \to -\ba$ と可換でないため、$f$, $g$ の段階では単純な規則にならない。

